% ===========================================================================
%  preambolo.tex — pacchetti, stile, macro e notazione
%  Motore previsto: pdfLaTeX + Biber.
% ===========================================================================

% --- Lingua e codifica -----------------------------------------------------
\usepackage[T1]{fontenc}
\usepackage[utf8]{inputenc}
% english è caricato perché il sommario contiene un abstract in inglese
% (\begin{otherlanguage}{english}); main=italian resta la lingua del documento.
\usepackage[english, main=italian]{babel}
\usepackage{lmodern}
\usepackage{microtype}
\usepackage{csquotes}          % richiesto da babel + biblatex

% --- Geometria di pagina ---------------------------------------------------
\usepackage[
  a4paper,
  inner=3.0cm, outer=2.5cm,
  top=2.8cm, bottom=2.8cm,
  headheight=15pt
]{geometry}
\usepackage{setspace}
\onehalfspacing

% --- Matematica ------------------------------------------------------------
\usepackage{amsmath,amssymb,amsthm,mathtools}
\usepackage{bm}
\usepackage{braket}            % \ket, \bra, \braket — notazione di Dirac

% --- Tabelle e figure ------------------------------------------------------
\usepackage{graphicx}
\usepackage{booktabs}
\usepackage{tabularx}
\usepackage{longtable}
\usepackage{multirow}
\usepackage{makecell}
\usepackage{rotating}
\usepackage{float}
\usepackage[font=small,labelfont=bf,width=0.92\textwidth]{caption}
\usepackage{subcaption}
\usepackage{siunitx}
\sisetup{
  output-decimal-marker = {,},
  group-separator = {\,},
  group-minimum-digits = 4,
  detect-weight = true,
  table-align-text-post = false
}

% --- Elenchi, note, intestazioni ------------------------------------------
\usepackage{enumitem}
\setlist{itemsep=2pt, topsep=4pt}
\usepackage{fancyhdr}
\pagestyle{fancy}
\fancyhf{}
\fancyhead[L]{\nouppercase{\small\leftmark}}
\fancyhead[R]{\small\thepage}
\renewcommand{\headrulewidth}{0.4pt}
\fancypagestyle{plain}{\fancyhf{}\fancyfoot[C]{\small\thepage}%
  \renewcommand{\headrulewidth}{0pt}}

% --- Codice sorgente -------------------------------------------------------
\usepackage{xcolor}
\definecolor{codebg}{gray}{0.96}
\definecolor{codekw}{RGB}{0,90,160}
\definecolor{codestr}{RGB}{160,60,20}
\definecolor{codecom}{RGB}{110,110,110}

\usepackage{listings}
\lstdefinestyle{python}{
  language=Python,
  basicstyle=\ttfamily\footnotesize,
  keywordstyle=\color{codekw}\bfseries,
  stringstyle=\color{codestr},
  commentstyle=\color{codecom}\itshape,
  backgroundcolor=\color{codebg},
  numbers=left, numberstyle=\tiny\color{codecom}, numbersep=8pt,
  frame=single, framerule=0pt, rulecolor=\color{codebg},
  showstringspaces=false,
  breaklines=true, breakatwhitespace=true,
  columns=fullflexible,
  captionpos=b,
  literate={à}{{\`a}}1 {è}{{\`e}}1 {é}{{\'e}}1 {ì}{{\`\i}}1 {ò}{{\`o}}1 {ù}{{\`u}}1
           {π}{{$\pi$}}1 {→}{{$\rightarrow$}}1 {·}{{$\cdot$}}1 {⟨}{{$\langle$}}1 {⟩}{{$\rangle$}}1
}
\lstdefinestyle{shell}{
  basicstyle=\ttfamily\footnotesize,
  backgroundcolor=\color{codebg},
  frame=single, framerule=0pt, rulecolor=\color{codebg},
  breaklines=true, breakatwhitespace=false,
  columns=fullflexible,
  literate={à}{{\`a}}1 {è}{{\`e}}1 {ì}{{\`\i}}1 {ò}{{\`o}}1 {ù}{{\`u}}1
}
\lstset{style=python}

% --- Annotazioni di lavorazione -------------------------------------------
% Disattivate per la consegna. Per riattivarle in fase di revisione togliere
% l'opzione [disable].
\usepackage[disable, textwidth=2.6cm, textsize=footnotesize, shadow]{todonotes}
\newcommand{\DA}[1]{\todo[color=orange!35]{#1}}          % dato da aggiornare
\newcommand{\VER}[1]{\todo[color=cyan!25]{#1}}           % da verificare sul codice
\newcommand{\SCR}[1]{\todo[inline, color=yellow!30]{#1}} % blocco da scrivere

% --- Bibliografia ----------------------------------------------------------
\usepackage[
  backend=biber,
  style=authoryear-comp,
  sorting=nyt,
  maxcitenames=2,
  maxbibnames=99,
  giveninits=true,
  uniquename=init,
  natbib=true,
  backref=false
]{biblatex}
\addbibresource{bibliografia.bib}

% --- Collegamenti ----------------------------------------------------------
\usepackage[hidelinks, pdfusetitle]{hyperref}
\hypersetup{
  pdftitle={\TitoloTesi{} --- \SottotitoloTesi},
  pdfauthor={\Candidato},
  pdfsubject={\TipoLaurea{} --- \CorsoDiLaurea},
  pdfkeywords={\ParoleChiave},
  bookmarksnumbered=true,
  linktoc=all
}
\usepackage[italian, nameinlink]{cleveref}
\crefname{table}{tabella}{tabelle}
\Crefname{table}{Tabella}{Tabelle}
\crefname{figure}{figura}{figure}
\Crefname{figure}{Figura}{Figure}
\crefname{equation}{equazione}{equazioni}
\Crefname{equation}{Equazione}{Equazioni}

% Ambienti semantici -------------------------------------------------------
\theoremstyle{definition}
\newtheorem{definizione}{Definizione}[chapter]
\newtheorem{ipotesi}{Ipotesi}[chapter]

\theoremstyle{remark}
\newtheorem*{implementazione}{Comportamento implementato}
\newtheorem*{osservazione}{Dato osservato}
\newtheorem*{interpretazione}{Interpretazione}
\newtheorem*{limite}{Limite}

% ===========================================================================
%  MACRO DI NOTAZIONE
%  Definite una volta sola: la coerenza notazionale fra i capitoli 3, 4, 5 e 7
%  dipende esclusivamente da queste definizioni.
% ===========================================================================

% Riferimenti a entità del codice ------------------------------------------
\newcommand{\code}[1]{\texttt{#1}}
\newcommand{\file}[1]{\texttt{\small #1}}
\newcommand{\classe}[1]{\texttt{#1}}

% Insiemi e operatori -------------------------------------------------------
\newcommand{\R}{\mathbb{R}}
\newcommand{\C}{\mathbb{C}}
\newcommand{\Hilb}{\mathcal{H}}
\newcommand{\E}{\mathbb{E}}
\DeclareMathOperator{\softmax}{softmax}
\DeclareMathOperator{\Tr}{Tr}
\DeclareMathOperator*{\argmax}{arg\,max}
\DeclareMathOperator{\arctanh}{arctanh}
\newcommand{\expval}[1]{\langle #1 \rangle}
\newcommand{\norm}[1]{\lVert #1 \rVert}
\newcommand{\abs}[1]{\lvert #1 \rvert}

% Forme tensoriali ----------------------------------------------------------
% \shape{B,T,C} produce (B, T, C) in stile matematico uniforme.
\newcommand{\shape}[1]{\ensuremath{(#1)}}

% Transformer (capitolo 3) --------------------------------------------------
\newcommand{\dbatch}{B}          % batch
\newcommand{\dseq}{T}            % lunghezza di sequenza (block_size)
\newcommand{\demb}{C}            % dimensione di embedding (n_embd)
\newcommand{\nhead}{H}           % numero di teste
\newcommand{\dhead}{d_h}         % dimensione per testa (head_size)
\newcommand{\dvoc}{V}            % dimensione del vocabolario

% Vision Transformer (capitolo 4) ------------------------------------------
\newcommand{\npatch}{N}          % numero di patch
\newcommand{\dpatch}{d_{\mathrm{p}}}  % dimensione della patch appiattita
\newcommand{\dvit}{d}            % dimensione di embedding del ViT (embed_dim)
\newcommand{\nch}{C_{\mathrm{ch}}}

% Circuito quantistico ------------------------------------------------------
\newcommand{\nq}{n_q}            % numero di qubit
\newcommand{\nqlayers}{L}        % profondità di StronglyEntanglingLayers
\newcommand{\Uenc}{U}

% Metriche e statistica -----------------------------------------------------
\newcommand{\gap}{G}                       % gap di generalizzazione
\newcommand{\acctr}{a_{\mathrm{train}}}
\newcommand{\accte}{a_{\mathrm{test}}}
\newcommand{\PPL}{\mathrm{PPL}}
\newcommand{\dz}{d_z}                      % dimensione dell'effetto appaiata
\newcommand{\dmean}{\bar{d}}
\newcommand{\sd}{s_d}
\newcommand{\IC}{\text{IC}_{95\%}}

% Convenzione di segno, richiamata nelle didascalie del capitolo 8 ----------
\newcommand{\convsegno}{%
  Convenzione di segno: le differenze sono calcolate come
  \emph{quantistico $-$ classico} (\code{direction: "quantum\_minus\_classical"}).}

% Etichette di esito delle ipotesi -----------------------------------------
\newcommand{\nonsupportata}{\textbf{non supportata}}
\newcommand{\evidenzacontraria}{\textbf{evidenza contraria}}
